\subsection{One primitive correlation and a one-sided target}
\label{ssec:primitive-energy}

The coprime coefficients in the preceding expansion also admit an
exponential weight without the local-energy integral. Put
\[
 C_\lambda(X)=\sum_{\substack{a,b\ge1\\(a,b)=1}}
       \lambda(a)\lambda(b)e^{-(a+b)/X}
       =\sum_{k\ge2}c(k)e^{-k/X}.
\]
The letter $C_\lambda$ here denotes this smoothed primitive correlation,
not a correlation at a fixed shift. It need not be nonnegative.

\begin{proposition}[Exact primitive decomposition]
\label{prop:primitive-decomposition}
For every $X>0$,
\begin{equation}\label{eq:primitive-decomposition}
 B(X)^2=\sum_{d\ge1}C_\lambda(X/d),\qquad
 C_\lambda(X)=\sum_{d\ge1}\mu(d)B(X/d)^2.
\end{equation}
Both series converge absolutely.
\end{proposition}
\begin{proof}
Partition the ordered pairs $(m,n)$ in $B(X)^2$ by their greatest
common divisor $d$, and write $m=da$, $n=db$. Complete
multiplicativity gives
$\lambda(da)\lambda(db)=\lambda(a)\lambda(b)$ because
$\lambda(d)^2=1$. This proves the first identity. To obtain the second,
insert
\[
 \mathbf1_{(m,n)=1}=\sum_{d\mid m,\ d\mid n}\mu(d)
\]
in the defining series for $C_\lambda$ and use the same identity for
the signs. All rearrangements are justified by the finite majorant
\[
 \sum_{d\ge1}\left(\sum_{a\ge1}e^{-da/X}\right)^2
       =\sum_{d\ge1}(e^{d/X}-1)^{-2}<\infty.
\]
For instance, its tail is bounded by a geometric series after
$d\ge X$. The same bound proves absolute convergence of the two
series in \eqref{eq:primitive-decomposition}.
\end{proof}

\begin{theorem}[A one-sided primitive criterion]
\label{thm:primitive-rh-criterion}
The following assertions are equivalent:
\begin{enumerate}
 \item RH;
 \item $B(X)\ll_\varepsilon X^{1/2+\varepsilon}$ for every
       $\varepsilon>0$ and $X\ge1$;
 \item for every $\varepsilon>0$ there exists $K_\varepsilon$ such that
       $C_\lambda(X)\le K_\varepsilon X^{1+\varepsilon}$ for every
       $X\ge1$.
\end{enumerate}
Only an upper bound is required in the third assertion.
\end{theorem}
\begin{proof}
Proposition~\ref{prop:smoothed} gives the equivalence of the first two
assertions. For the implication from the third to the second, first
enlarge $K_\varepsilon$ so that its bound holds for all $X>0$. This
is possible because
\[
 |C_\lambda(X)|\le(e^{1/X}-1)^{-2}
\]
decays faster than any fixed power as $X\downarrow0$. The first
identity in \eqref{eq:primitive-decomposition} now yields
\begin{equation}\label{eq:primitive-upper-transfer}
 0\le B(X)^2\le K_\varepsilon\zeta(1+\varepsilon)X^{1+\varepsilon}.
\end{equation}
Taking $\varepsilon=2\delta$ proves the second assertion with exponent
$1/2+\delta$, for every $\delta>0$.

Conversely, use the second assertion with exponent
$1/2+\varepsilon/2$, and extend the resulting bound to $0<X<1$
by exponential decay. Write the squared bound as
$|B(X)|^2\le A_\varepsilon X^{1+\varepsilon}$. The second identity
in \eqref{eq:primitive-decomposition} then gives
\begin{equation}\label{eq:primitive-absolute-converse}
 |C_\lambda(X)|\le
 A_\varepsilon\frac{\zeta(1+\varepsilon)}{\zeta(2+2\varepsilon)}
 X^{1+\varepsilon}.
\end{equation}
Here the squarefree Euler product gives
$\sum_{d\ge1}|\mu(d)|d^{-1-\varepsilon}
=\zeta(1+\varepsilon)/\zeta(2+2\varepsilon)$, with absolute
convergence. This proves the third assertion.
\end{proof}

The theorem isolates a single weighted upper estimate. No bound for
each individual $c(k)$ or each fixed additive shift is needed.
However, the upper estimate for $C_\lambda$ has not been proved.
The exact identities do not supply it, and positivity of $B(X)^2$
does not bound $C_\lambda$ from above.
\footnote{The constant and all sufficiently large scales are part of
the third assertion. Checking finitely many coefficients or scales
does not establish it.}

For exact finite controls, put
\[
 A(k)=\sum_{a=1}^{k-1}\lambda(a)\lambda(k-a),\qquad A(1)=c(1)=0.
\]
Coefficientwise the same proof gives
\begin{equation}\label{eq:primitive-finite-controls}
 A(k)=\sum_{d\mid k}c(k/d),\qquad
 c(k)=\sum_{d\mid k}\mu(d)A(k/d).
\end{equation}
The checker \texttt{anc/check\_primitive\_identities.py} compares
these divisor identities with direct greatest-common-divisor
enumeration for $1\le k\le1024$. Its scope is finite identities,
not the asymptotic assertion in Theorem~\ref{thm:primitive-rh-criterion}.

\begin{figure}[tb]
\centering
\begin{tikzpicture}[>=Latex,node distance=1.15cm,font=\small]
 \node[draw,rounded corners,fill=orange!10,align=center,
       text width=8.1cm] (input) at (0,0)
   {unproved input for the actual Liouville function\\
    $C_\lambda(X)\le K_\varepsilon X^{1+\varepsilon}$ at every scale};
 \node[draw,rounded corners,fill=blue!6,align=center,
       text width=8.1cm] (sum) at (0,-1.4)
   {$B(X)^2=\sum_{d\ge1}C_\lambda(X/d)$\\
    $B(X)^2\le K_\varepsilon\zeta(1+\varepsilon)X^{1+\varepsilon}$};
 \node[draw,rounded corners,fill=blue!6,align=center,
       text width=8.1cm] (zero) at (0,-2.8)
   {Mellin holomorphy in $\R s>1/2$\\
    and exclusion of right-hand zeta zeros};
 \draw[->,thick] (input)--(sum);
 \draw[->,thick] (sum)--(zero);
\end{tikzpicture}
\caption{Proved implications from an unproved upper estimate. The
primitive decomposition retains the signs; it does not prove its input.}
\label{fig:primitive-criterion}
\end{figure}
